\section{NUMA effects}
\label{section:node-level:numa}

%Gibt es ein Amdahl law?
%   \item[$\boxplus$] Use Amdahl's law to characterise your code. If it fits to
%   the law, i.e.~there is an $f_{\text{intra}}$ such that
%   \eqref{equation:amdahl:speedup}.

% Allocate memory where you need it
% APS can provide this indicator

On modern cluster nodes, with large numbers of CPU cores, the shared memory is
often split into distinct regions known as Non-Unified Memory Access (NUMA) or
cache-coherent NUMA (ccNUMA) domains. These behave as shared memory systems,
however, because they are not the same physical block of memory some accesses
have to pass through an interconnect rather than the through the cache
hierarchy. This incurs a performance penalty for memory accesses across NUMA
domains. In short, we want to ensure that, as much as possible, memory accesses
are carried within NUMA domains, rather than across them.

\subsection{NUMA tools and performance counters}

There are multiple ways of detecting NUMA issues but we begin with a simple
test which can imply NUMA effects are present. The \texttt{numactl} command is
common on Linux commands and is responsible for controlling NUMA policies on
systems, including how memory is laid out across processes. Without digging
into the details of \texttt{numactl} at this stage, one simple test for NUMA
effects is listed below.

\begin{instructions}{\texttt{numactl}}
  We can run our code executable, \texttt{code}, through the \texttt{numactl} command 
  \begin{verbatim}
    numactl --interleave ./code
  \end{verbatim}
  If there is a difference in the runtime of the code with and without this prepended NUMA command, then this can imply NUMA effects are present.
\end{instructions}

If our \texttt{numactl} test implies that NUMA effects may be present, then we
can explore further by using performance counters. 

\begin{instructions}{\likwid}
  We can once again turn to \likwid to run our executable, \texttt{code}, this time accessing NUMA performance counters with the \texttt{NUMA} performance group 
  \begin{verbatim}
    likwid-perfctr -g NUMA -c 0-3 ./code
  \end{verbatim}
  Where we have selected a pinning of the first four cores but this is just illustrative.
\end{instructions}

NUMA-related performance counters and metrics can vary a lot across systems, so
these diagnostics should all be used to build up a picture, rather than being
direct evidence of specific problems.

\subsection{First-touch policy}

Often NUMA effects stem from the lack of the so-called `first-touch policy', in
which memory is allocated close to the processes that first `touched' it. Here,
by `touch' we mean \emph{write to}. This could, for example, be during the
initialisation stage of a code in which objects are being created. If these
initialisations are done in serial, but the code is OpenMP parallel, the memory
pages for these allocations will be mapped to the process running the
initialisations. This will then create a performance penalty during the runtime.

We do not here work through how to engineer codes to avoid NUMA effects but, in
short, when initialising data for a parallel application it is best to
initialise within a parallel region to improve data locality. For example, when
creating an array in an OpenMP application, wrapping the initialisation in an
OpenMP parallel pragma will assign the data to the local memory of the parallel
processes. Likewise, we need to ensure that the same OpenMP schedule is used in
the initialisation as in all other computational loops, and this can be
achieved by using a \emph{static schedule}. We can also use \texttt{numactl} to
allocate memory in a round-robin fashion. Similarly, on many systems NUMA
balancing maybe switched on though again it is typically good practice to
implement a first-touch policy

This can help us assess the presence NUMA effects because we can focus our
analysis on initialisation steps to see if a first-touch policy has been
included. As in previous sections, if we want to be kernel specific in our
analysis we can use the \likwid marker API to wrap loops and kernels to access
performance counters to get further details on the NUMA effects on kernels that
contribute significantly to runtime.
